% !TeX program = pdflatex
%
% "Modular anchors for Apery-like companions" -- revised per the hostile
% referee report of 2026-08-05 (work/SOL_REVIEW_MODULAR_ANCHORS.md, with the
% point-by-point response appended there). Working paper of River's odd-zeta
% programme.
%
% Provenance: work/MODULAR_COACTION_PROBE.md, work/SPORADIC_MODULAR_DICTIONARY.md,
% work/Z5_MODULARITY_PROBE.md (incl. 2026-08-05 correction),
% work/MODULAR_FACTORY.md, work/CURVE_BLINDNESS.md, work/DELTA_MODULAR_THEOREM.md,
% work/z5eps/eps40--eps57 with logs.  Companion papers:
% papers_out/variational_anchor, papers_out/expository_apery, papers_out/half_apery.
%
% Evidence labels are load-bearing and are never silently upgraded.
%
\documentclass[11pt,reqno]{amsart}

\usepackage[T1]{fontenc}
\usepackage[utf8]{inputenc}
\usepackage{amsmath,amssymb,amsthm}
\usepackage{mathtools}
\usepackage{booktabs}
\usepackage{enumitem}
\usepackage[margin=1.05in]{geometry}
\usepackage{microtype}
\usepackage[hidelinks]{hyperref}

\theoremstyle{plain}
\newtheorem{theorem}{Theorem}[section]
\newtheorem{proposition}[theorem]{Proposition}
\newtheorem{lemma}[theorem]{Lemma}
\newtheorem{corollary}[theorem]{Corollary}
\newtheorem{conjecture}[theorem]{Conjecture}

\theoremstyle{definition}
\newtheorem{definition}[theorem]{Definition}
\newtheorem{census}[theorem]{Computer-assisted census}
\newtheorem{verification}[theorem]{Computer-assisted finite verification}
\newtheorem{problem}[theorem]{Problem}

\theoremstyle{remark}
\newtheorem{remark}[theorem]{Remark}
\newtheorem{observation}[theorem]{Observation}

\newcommand{\Z}{\mathbb{Z}}
\newcommand{\Q}{\mathbb{Q}}
\newcommand{\bin}[2]{\binom{#1}{#2}}
\newcommand{\HH}[2]{H^{(#1)}_{#2}}
\newcommand{\Lbz}{L_{\mathrm{BZ}}}
\newcommand{\thq}{\theta_q}
\newcommand{\tht}{\theta_t}

% evidence-class markers
\newcommand{\Thm}{\textup{\textsf{[THM]}}}
\newcommand{\Proved}{\textup{\textsf{[PROVED]}}}
\newcommand{\Verified}{\textup{\textsf{[VERIFIED]}}}
\newcommand{\VerifiedR}[1]{\textup{\textsf{[VERIFIED: #1]}}}
\newcommand{\Excluded}{\textup{\textsf{[EXCLUDED]}}}
\newcommand{\ExcludedR}[1]{\textup{\textsf{[EXCLUDED: #1]}}}
\newcommand{\Open}{\textup{\textsf{[OPEN]}}}

\begin{document}

\title[Modular anchors for Ap\'ery-like companions]
      {Modular anchors for Ap\'ery-like companions:\\
       second solutions as Eichler integrals,\\
       with the symmetric-square theorem for the sporadic operators}

\author{}
\date{\today}

\begin{abstract}
For an Ap\'ery-like sequence $A(n)$ --- one of the fifteen sporadic pairs, or
a row of the Brown--Zudilin cellular family for $\zeta(5)$ --- the
\emph{companion problem} asks for a closed formula for the second solution
$B(n)$ of its recurrence. We prove an abstract \emph{companion inversion
theorem}: if the generating operator is exactly rectified by its nome,
$L=C(t)\,F\,\thq^{r}F^{-1}$, then the normalized companion is the
Eichler-type integral $y_B=F\,\thq^{-r}\bigl(t/(C F)\bigr)$, for all $n$.
We then prove the rectification hypothesis for all fifteen sporadic
families: for the nine second-order families it is Frobenius theory; for
the six third-order families we prove, by exact operator identities over
$\Q(t)$, that each is the symmetric square of an explicit second-order
operator $\widetilde D=\theta^2-t(2a\theta^2+a\theta+\tfrac b2)
+ct^2(\theta+\tfrac12)^2$ --- a uniform closed form. Consequently the
companion formula holds for all $n$ for all fifteen families, including the
seven for which we show (computer-assisted, precisely scoped) that no
companion exists in the complete tame harmonic ansatz of the programme.
Supporting results: the modular identification
$F_\zeta=\eta_3^{10}/(\eta_1\eta_9)^3=\tfrac18(9E_2(q^9)-E_2(q))$ is proved
(Ligozat--Sturm); a first annihilation law for the family $\zeta$'s summand
is proved by a WZ certificate; a new uniform weight-one congruence
$A(p)\equiv b_p-b\,\chi_N(p)\pmod p$ survives sixteen held-out primes per
family; a curve-blindness theorem for the $\zeta(5)$ top row is stated with
a degree-independent reduction (coupling lemma plus a proved Veronese
saturation) and finite certified residual computations; and a
census of $\sim5\times10^4$ modular pairs finds no sixteenth sporadic pair
in its precisely defined search universe. Mirror-type local and arithmetic
signatures of the $\zeta(5)$ operator (integral mirror map, integral
Yukawa-type datum) and a coaction proposal are collected in a clearly
labelled conjectural section. Every statement carries its evidence label;
finite computations are never reported as proofs.
\end{abstract}

\maketitle
\setcounter{tocdepth}{1}
\tableofcontents

%======================================================================
\section{Introduction}\label{sec:intro}

\subsection{The companion problem, and what is proved here}
Ap\'ery's proof of the irrationality of $\zeta(3)$ rests on a pair: the
integral solution $a_n=\sum_k\bin nk^2\bin{n+k}k^2$ of his recurrence and
the second solution $b_n$ with $b_n/a_n\to\zeta(3)$. Companion formulas ---
closed expressions for second solutions --- are the engine of the
denominator and congruence arithmetic in the surrounding programme. The
classical shape is a binomial sum against a harmonic weight; for seven of
the fifteen sporadic Ap\'ery-like families no such formula was known, and
the programme's searches (\S\ref{sec:negatives}) indicate none exists in
the natural finite ansatz.

This paper's main line is positive and, in its core, proved:

\begin{enumerate}[leftmargin=2em]
\item an abstract \emph{companion inversion theorem}
(Theorem~\ref{thm:inversion}): exact rectification of the operator by its
nome implies a closed Eichler-integral formula for the companion, valid for
all $n$;
\item the \emph{rectification theorem for the fifteen}
(Theorem~\ref{thm:sym2}): the nine second-order sporadic operators are
rectified by Frobenius theory, and each of the six third-order operators is
proved to be the symmetric square of an explicit second-order operator, by
exact identities in $\Q(t)$;
\item hence the \emph{uniform companion theorem}
(Theorem~\ref{thm:companion}): for every one of the fifteen sporadic pairs,
\[
B(n)\;=\;[t^n]\;F\cdot\thq^{-r}\!\Bigl(\frac{t}{C(t)\,F}\Bigr)
\qquad\text{for all }n\ge0 ,
\]
with $r$ the operator order, $q$ the nome, $F=y_0(t(q))$, and
$C=P_r/\sigma^r$ the explicit rectification factor.
\end{enumerate}

Around this proved core we organise, in separate tiers with separate
epistemic labels: exact modular results (\S\ref{sec:modular}), precisely
scoped computer-assisted negatives (\S\ref{sec:negatives}), arithmetic
congruence laws (\S\ref{sec:cong}), and a conjectural mirror/coaction
programme for the $\zeta(5)$ family (\S\ref{sec:mirror}).

\subsection{Notational dictionary}\label{sec:dict}
Throughout, $r$ denotes the \emph{order of the generating differential
operator} $=$ the order of the three-term recurrence's companion matrix
$+1$ in the classical normalisation: $r=2$ for the nine ``weight-2''
sporadic rows, $r=3$ for the six ``weight-3'' rows. For the families with
\emph{identified} parametrizations the modular weight of $F$ is $r-1$, and
the congruence weight $w$ of the programme's tables equals $r$; for the
three unidentified Cooper rows ($r=3$) no modular weight is asserted, and
statements grouping them by weight refer to the congruence weight $w=r$
only. We use $r$ exclusively in operator
statements; earlier drafts' floating ``$w$'' has been retired
(referee report, \S2). $\theta_x=x\,d/dx$; $t$ is the recurrence variable,
$q$ the nome; $\sigma:=\tht\log q\in1+t\Q[[t]]$.

\subsection{Evidence discipline}\label{sec:labels}
\Thm/\Proved: complete proof (human-readable; computer algebra used only
for exact rational identities, which are finite and displayed or scripted).
\Verified: exact rational or modular computation over a stated finite
range. \Excluded: a negative established either by proof or by an exact
finite computation over a precisely defined search space, with the space
stated. \Open: open. Two standing rules, adopted at the referee's
insistence: finite coefficient agreement is never called an identity, and
finite integrality is never called modularity or arithmeticity
(\S\ref{sec:modular}, \S\ref{sec:mirror}). A point-by-point response to
the referee report, and a log guarding against silent evidential
upgrades in future revisions, is maintained at
\texttt{work/SOL\_REVIEW\_MODULAR\_ANCHORS.md}.

\begin{lemma}[modular-to-rational transfer]\label{lem:modrat}
\Proved\
Let $Mx=v$ be a linear system over $\Q$, $p$ a prime dividing no
denominator of the entries of $(M,v)$, and suppose the reduction
$\bar M\bar x=\bar v$ over $\mathbb F_p$ is inconsistent. Then $Mx=v$ is
inconsistent over $\Q$.
\end{lemma}

\begin{proof}
A rational solution $x$ has some denominator $d$; choosing $p\nmid d$
would reduce it to a solution mod $p$. In general, if $x\in\Q^m$ solves
the system, then for any prime $p$ not dividing the denominators of
$M,v,x$, the reduction solves the reduced system. Inconsistency mod $p$
therefore excludes all rational solutions whose denominators avoid $p$;
running the computation at two unrelated primes excludes rational
solutions except those whose denominators are divisible by both, and the
certificate of inconsistency (a left null vector $u$ with $u^TM=0$,
$u^Tv\ne0$ over $\mathbb F_p$) lifts: since the entries of $M,v$ are
$p$-integral, any rational solution reduces mod $p$ after clearing its
denominator, contradicting $u^T\bar v\neq0$ unless the denominator is
divisible by $p$. As no rational vector has denominator divisible by
every good prime, inconsistency at a single good prime already implies
rational inconsistency; we use two primes as an independent guard against
implementation error, not as a logical necessity.
\end{proof}

All \Excluded\ statements below are of this form: the matrices are exact
reductions of rational systems (integer numerators, denominators coprime
to the $22$-bit primes $4194301$, $4194247$), so the negatives are
theorems over $\Q$ about the stated finite ansatz spaces.

%======================================================================
\section{The companion inversion theorem}\label{sec:core}

\begin{definition}[nome and rectification]\label{def:nome}
Let $L$ be an order-$r$ operator in $\tht$ over $\Q(t)$ whose solution
space at $t=0$ has a Frobenius basis $y_0=1+O(t)$,
$y_1=y_0\log t+g$, \dots. The \emph{nome} is $q=t\exp(g/y_0)$, so
$\log q=y_1/y_0$; write $t(q)$ for the inverse and $F(q)=y_0(t(q))$. We
say $L$ is \emph{exactly rectified} if
\begin{equation}\label{eq:rect}
L\;=\;C(t)\cdot F\cdot\thq^{\,r}\cdot F^{-1}
\end{equation}
for a rational function $C$; equivalently, the kernel of $L$ is
$\operatorname{span}\{F,F\log q,\dots,F(\log q)^{r-1}/(r-1)!\}$, i.e.\ the
full Frobenius basis collapses onto powers of the single logarithm
$\log q$. For $r\ge3$ this is strictly stronger than maximal unipotent
monodromy (referee report, \S1; see Theorem~\ref{thm:sym2}).
\end{definition}

\begin{lemma}[boundary defect]\label{lem:defect}
\Proved\
Let $L$ come from an R2 or R3 sporadic recurrence and let $y_B=\sum
B(n)t^n$ with $B(0)=0$, $B(1)=1$ the normalized second solution. Then
$L(y_B)=t$ exactly: the coefficients of $t^m$, $m\neq1$, vanish because
they are instances of the recurrence, and the coefficient of $t$ is
$1^{\,r}B(1)-(\text{lower terms})\cdot B(0)=1$.
\end{lemma}

\begin{lemma}[uniqueness]\label{lem:unique}
\Proved\
A log-free power series $y$ with $y(0)=0$ and $L(y)=t$ is unique: two such
differ by a log-free kernel element vanishing at $0$; the log-free part of
the kernel is $\Q\,y_0$ under \eqref{eq:rect} (all other basis elements
carry logarithms), and $y_0(0)=1\neq0$.
\end{lemma}

\begin{theorem}[companion inversion]\label{thm:inversion}
\Thm\
Suppose $L$ is exactly rectified as in \eqref{eq:rect}. Then
\[
y_B\;=\;F\cdot\thq^{-r}\!\Bigl(\frac{t}{C(t)\,F}\Bigr),
\]
where $\thq^{-1}$ is the constant-free antiderivative (well defined since
$t/(CF)\in q\,\Q[[q]]$), and consequently
$B(n)=[t^n]\,F\,\thq^{-r}(t/(CF))$ for all $n\ge0$.
\end{theorem}

\begin{proof}
The displayed series is log-free, vanishes at $t=0$, and satisfies
$L(\cdot)=C\,F\,\thq^r(F^{-1}\cdot)=C\,F\cdot\bigl(t/(CF)\bigr)/1=t$ by
\eqref{eq:rect}. Lemmas \ref{lem:defect} and \ref{lem:unique} identify it
with $y_B$. (Intermediate antiderivative constants alter the result by
kernel elements with logarithms or by $\lambda F$; both are excluded by
log-freeness and $y(0)=0$.)
\end{proof}

%======================================================================
\section{Rectification for the fifteen: the symmetric-square theorem}
\label{sec:sym2}

\subsection{Order two}
\begin{lemma}\label{lem:r2}
\Proved\
Every order-$2$ operator with Frobenius basis $y_0$, $y_1=y_0\log t+g$ is
exactly rectified: in the $q$-coordinate its kernel is
$\{F,F\log q\}$ by definition of $q$, which is the kernel of
$F\thq^2F^{-1}$; comparing leading coefficients gives
\eqref{eq:rect} with $C=P_2/\sigma^2$, $P_2$ the $\tht^2$-coefficient of
$L$. This covers the nine second-order sporadic families
($\mathbf A$--$\mathbf F$, $s_7$, $s_{10}$, $s_{18}$), including the
$d\neq0$ Cooper rows.
\end{lemma}

\subsection{Order three}
For an R3 family $(a,b,c)$ (all six have $d=0$) the generating operator is
\[
L=\theta^3-t(2\theta+1)(a\theta^2+a\theta+b)+c\,t^2(\theta+1)^3,
\qquad
P_3(t)=1-2at+ct^2 .
\]

\begin{theorem}[the six symmetric squares]\label{thm:sym2}
\Proved\ (exact identities in $\Q(t)$;
\texttt{work/z5eps/eps57\_sym2\_all.py})\quad
For each of the six third-order sporadic families, $L$ equals
$t^3P_3(t)$ times the symmetric square of the second-order operator
\[
\widetilde D\;=\;\theta^2-t\Bigl(2a\theta^2+a\theta+\tfrac b2\Bigr)
+c\,t^2\Bigl(\theta+\tfrac12\Bigr)^2 ,
\]
uniformly in the family parameters: writing $L$ in $\partial$-form and
matching against
$\mathrm{Sym}^2(\partial^2+p\partial+r_\ast)=\partial^3+3p\partial^2
+(2p^2+p'+4r_\ast)\partial+(4pr_\ast+2r_\ast')$, the forced
$(p,r_\ast)$ agree with $\widetilde D$'s and the residual
$A_0/A_3-(4pr_\ast+2r_\ast')$ is \emph{identically zero} in $\Q(t)$, for
$\alpha, \gamma, \delta, \varepsilon, \zeta, \eta$ separately. For
example, for $\gamma$ (Ap\'ery $\zeta(3)$),
$p=(2t^2-51t+1)/\bigl(t(t^2-34t+1)\bigr)$,
$r_\ast=(t-10)/\bigl(4t(t^2-34t+1)\bigr)$.
\end{theorem}

\begin{corollary}[rectification for all fifteen]\label{cor:rect15}
\Proved\
Each of the fifteen sporadic operators is exactly rectified: for $r=2$ by
Lemma~\ref{lem:r2}; for $r=3$ because, with $u,v$ a Frobenius basis of
$\widetilde D$, the kernel of $L$ is $\{u^2,uv,v^2/2\}$ and
$v/u=\log q$ (the nome of $\widetilde D$ equals the nome of $L$ since
$y_1/y_0=uv/u^2$), so the kernel is
$\{F,F\log q,\tfrac12F\log^2q\}$; the leading-coefficient comparison gives
$C=P_3/\sigma^3$.
\end{corollary}

\begin{theorem}[uniform companion formula, all $n$]\label{thm:companion}
\Thm\
For every one of the fifteen sporadic pairs, with $r$ the operator order,
\[
B(n)\;=\;[t^n]\;F\cdot\thq^{-r}\!\Bigl(\frac{t\,\sigma^r}{P_r\,F}\Bigr)
\qquad\text{for all }n\ge0 .
\]
\end{theorem}

\begin{proof}
Theorem~\ref{thm:inversion} with Corollary~\ref{cor:rect15} and
$C=P_r/\sigma^r$.
\end{proof}

\begin{verification}\label{ver:companion}
\VerifiedR{exact $\Q$, $n\le20$ per family ($n\le22$ for $\zeta$,
$n\le26$ for $\gamma$); \texttt{eps52}}\ The formula of
Theorem~\ref{thm:companion} was independently confirmed coefficientwise
for all fifteen families; with the theorem now proved, these computations
serve as implementation checks.
\end{verification}

\begin{remark}
The referee's structural point stands recorded in
Definition~\ref{def:nome}: rectification at $r\ge3$ is strictly stronger
than MUM, and a generic MUM operator retains a Yukawa-type coupling in
canonical coordinates. What Theorem~\ref{thm:sym2} shows is that the
sporadic zoo sits entirely in the rectifiable class --- order-appropriate
in each case (referee report \S6) --- while the $\zeta(5)$ five-block does
not (\S\ref{sec:mirror}). The half-integral shift $\theta+\tfrac12$ in
$\widetilde D$ is the source of the ``half-Ap\'ery'' phenomena studied in
the companion paper \texttt{papers\_out/half\_apery}.
\end{remark}

%======================================================================
\section{Exact modular results}\label{sec:modular}

\subsection{The $\zeta$ identification is an identity}
\begin{proposition}\label{prop:zetaid}
\Proved\ (Ligozat conditions $+$ Sturm bound; \texttt{eps57}, Part 2)\quad
The eta quotient $\eta_3^{10}/(\eta_1\eta_9)^3$ is a holomorphic modular
form of weight $2$, level $9$, trivial character (cusp orders $0,1,0$ at
the cusps $1,\,1/3,\,\infty$ of $\Gamma_0(9)$; $\prod d^{e_d}=81$ a
square), and equals $\tfrac18\bigl(9E_2(q^9)-E_2(q)\bigr)$ in
$M_2(\Gamma_0(9))$: the Sturm bound is $2$ and agreement holds to
$q^{40}$.
\end{proposition}

\subsection{The first proved annihilation law for $\zeta$'s letters}
\begin{proposition}\label{prop:ann}
\Proved\ (order-$2$ Zeilberger certificate with symbolically verified
certificate identities, endpoint vanishing, and upward induction;
\texttt{work/z5eps/eps45\_zeta\_kernel.md}, \texttt{eps45\_verify1.py})\quad
For the summand $S_\zeta(n,k,l)=\bin nk^2\bin nl\bin kl\bin{k+l}n$ and all
$0\le k\le n$,
\[
\sum_{l}S_\zeta(n,k,l)\,\bigl(\HH1k-2\HH1l+\HH1{n-l}\bigr)\;=\;0 .
\]
Order $2$ is minimal for the certificate (order $1$ excluded by
residual-checked exhaustion). A second, weight-two annihilation law
remains \VerifiedR{exact $\Q$, all cells $n\le25$} with a documented
obstruction record (Gosper and joint order-$2$--$4$ solves excluded, also
over an expanded alphabet at orders $2$--$3$) \Open. Under the inner sum
the letters $\HH1l$, $\HH1{n-l}$ are therefore provably not free ---
every previous fit for $\zeta$ was conducted in a degenerate alphabet.
\end{proposition}

\subsection{Identifications at discovery grade}
\begin{verification}[the dictionary]\label{ver:dict}
\VerifiedR{coefficientwise to $q^{25}$ with zero-tail checks;
\texttt{eps51}; table in \texttt{work/SPORADIC\_MODULAR\_DICTIONARY.md}}\quad
All fifteen nome expansions are integral at scale one to the stated
order, and twelve of fifteen $(t,F)$ pairs agree to $q^{25}$ with explicit
(generalized-)eta quotients --- including
$\delta$ (level $12$: $t=q\,\eta_1^4\eta_4^4\eta_6^{16}/
(\eta_2^{16}\eta_3^4\eta_{12}^4)$,
$F=\eta_2^{12}\eta_3\eta_{12}/(\eta_1^3\eta_4^3\eta_6^4)$), apparently the
first modular data attached to $\delta$;
$\eta$ (level $20$); $\mathbf B$ (level $36$, with the half-integral
$q$-shift anomaly $\sum m\,e_m=132\notin24\Z$ reflecting the $b(-q)$
twist, and an equivalent untwisted level-$9$ representative found by the
census of \S\ref{sec:factory}). Cooper's three are integral but not
eta-type in this coordinate \Open.

Finite integrality and finite agreement are \emph{discovery} statements:
they do not by themselves assert modular parametrisation (referee report
\S4). The upgrade path is exact operator verification;
Proposition~\ref{prop:zetaid} completes the form-identification step for
$\zeta$, and for $\delta$ the parallel programme is carried out in
\texttt{work/DELTA\_MODULAR\_THEOREM.md}: Ligozat holomorphy of both
quotients is unconditional, the eta-side annihilator within the order-$3$
ansatz is exactly one-dimensional and equals $\delta$'s operator
coefficientwise, and one hypothesis remains (the degree-one property of
$t_\delta$ on the correct Atkin--Lehner extension) \Open.
\end{verification}

%======================================================================
\section{Computer-assisted negatives, precisely scoped}\label{sec:negatives}

All negatives in this section are of the shape of Lemma~\ref{lem:modrat}:
rational inconsistency of a fully specified finite linear system,
certified at two unrelated $22$-bit primes with identical ranks.

\subsection{The tame harmonic ansatz}\label{sec:exhaust}
\begin{definition}[the ansatz]\label{def:ansatz}
For a family with summand $S$ and congruence weight $w$, the \emph{complete
tame degree-$\le3$ ansatz} is the finite $\Q$-vector space of weights
$\mathfrak w=\sum_jc_jM_j$ where $M_j$ ranges over monomials of total
harmonic weight exactly $w$, of degree $\le3$ in weight-one letters, in
the letters $\HH sx$ (and, in the twisted variant, $K^{(s)}_\chi(x)=
\sum_{j\le x}\chi(j)j^{-s}$ with $\chi$ the conductor-matched character),
with arguments $x$ ranging over the family's complete tame argument set:
for $\zeta$: $\{n,k,n{-}k,l,n{-}l,k{-}l,k{+}l{-}n\}$;
for $\eta$: $\{n,k,2k,3k,4k,5k,n{-}k,n{-}2k,n{-}3k,n{-}4k,n{-}5k\}$;
for $\mathbf F$: $\{n,k,n{-}k,l,k{-}l,n{-}l,n{-}k{+}l\}$.
The companion equation is the linear system
$B(n)=\sum_{\text{cells}}S\,\mathfrak w$ over the stated fitting rows.
\end{definition}

\begin{theorem}[exhaustion of the ansatz]\label{thm:exhaust}
\ExcludedR{two primes, identical ranks; excess $90$--$173$;
\texttt{eps44}}\quad
For $\mathbf F$, $\zeta$, $\eta$, the system of
Definition~\ref{def:ansatz} is inconsistent over $\Q$, pure and twisted
(columns/rank: $35/23$, $119/75$; $140/125$, $770/734$; $418/418$,
$2530/2530$). Together with the programme's prior exclusions this
exhausts the complete tame degree-$\le3$ ansatz for all seven conjectural
families. The positive control (Franel) re-derives its known companion
exactly in the same harness with held-out validation.
\end{theorem}

\begin{remark}[scope]
This excludes exactly the stated ansatz. It does not exclude: higher
degree in weight-one letters, nested harmonic sums, rational-function
coefficients $c_j(n,k)$, other hypergeometric representations of the same
$A(n)$, or telescoping/indefinite-sum extensions (referee report \S3).
The claim ``no harmonic companion exists'' does not appear in this paper;
Theorem~\ref{thm:companion} makes the point that matters --- the companion
\emph{does} exist, in modular letters, for all fifteen --- and
Proposition~\ref{prop:ann} explains part of the ansatz failure
(degenerate alphabets).
\end{remark}

\subsection{Curve-blindness for the $\zeta(5)$ top row}\label{sec:blind}
\begin{definition}[curve atoms, pinning, boxes]\label{def:atoms}
Write the Brown--Zudilin cell
$T=\prod_L(x_L!)^{p_L}$ over the nine letters
$x_L\in\{n,k,l,n{\pm}k,n{\pm}l,k{+}l,n{+}k{+}l\}$,
$p_L=(1,-3,-3,1,1,-2,-2,-1,1)$. A \emph{curve atom of degree $D$} is
$u=(u_1,\dots,u_D)$, $u_i\in\Z^3$, denoting the $\Gamma$-continued cell
$T(n+\alpha(\varepsilon),k+\beta(\varepsilon),l+\gamma(\varepsilon))$,
$(\alpha,\beta,\gamma)=\sum_iu_i\varepsilon^i$; per letter the shift is
$s_L=\sum_id_{i,L}\varepsilon^i$ with $d_{i,L}$ the letter's linear form
on $u_i$. \emph{Pinning} of a $\Q$-combination means: the
$\varepsilon^1$ and $\varepsilon^2$ rows vanish and the $\varepsilon^3$
row lies in $\Q Q\oplus\Q\hat P$, imposed separately in every
$\zeta$-graded component (variant A adds the analogous $\varepsilon^4$
condition). The \emph{boxes} are $u_1\in\{-2..2\}^3$,
$u_2\in\{-1,0,1\}^3$, $u_3\in\{0,\pm e_i\}$.
\end{definition}

\begin{lemma}[coupling; degree-independent]\label{lem:coupling}
\Proved\
For any polynomial curve atom of any degree, the top-graded (harmonic
weight $=$ $\varepsilon$-order $r$) content of $[\varepsilon^r]$ depends
only on the first-order direction $u_1$, through the pure powers
$d_1^{\otimes r}$; all higher curve coefficients feed only lower weights.
Hence the top-grade question for \emph{all} polynomial curves reduces to
the span of the degree-$r$ Veronese images of first-order directions.
\end{lemma}

\begin{lemma}[Veronese saturation]\label{lem:sat}
\Proved\ (exact rank computation; \texttt{eps57\_heldout.py}, Part B)\quad
The directions of the box $\{-2..2\}^3$ have degree-$d$ Veronese images of
full rank $\dim\mathrm{Sym}^d(\Q^3)$ for $d=2,3,4,5$ ($6,10,15,21$). Hence
the finite computations below certify the top-graded component for
\emph{every} first-order direction in $\Z^3$, and so, by
Lemma~\ref{lem:coupling}, for every polynomial curve of every degree.
\end{lemma}

\begin{theorem}[curve-blindness]\label{thm:blind}
The top-graded $\varepsilon^5$ content of any pinned $\Q$-combination of
polynomial-curve atoms (any degree, any directions) has vanishing
$P_n$-component; and on the boxes of Definition~\ref{def:atoms} the full
statement holds: the entire $\varepsilon^5$ row of a pinned combination is
confined to $\Q Q\oplus\Q\hat P$, in the rational part and every
$\zeta$-graded component
\ExcludedR{two primes, identical ranks $53/36$ (linear), $88/59$
(quadratic), $80/54$ (cubic); \texttt{eps41}, \texttt{eps42}}. The
$\varepsilon^3$ coefficient meanwhile realises $\hat P$ freely. The
top-grade clause is unconditional by Lemmas
\ref{lem:coupling}--\ref{lem:sat}; the sub-top clauses are certified on
the stated boxes, where the pinning constraints (which involve $u_2,u_3$)
live.
\end{theorem}

\subsection{The census}\label{sec:factory}
\begin{definition}[search universe]\label{def:universe}
The census banks are: hauptmodul candidates
$t=q\prod_m(1-q^{mk})^{e_m}$ with $\sum e_m=0$, $\sum m\,e_m=24$; form
candidates $F=\prod_m(1-q^{mk})^{e_m}$ with $\sum e_m=2(r-1)$,
$\sum m\,e_m=0$, $r-1\in\{1,2\}$; Eisenstein members
$(d_2E_2(q^{d_2})-d_1E_2(q^{d_1}))/(d_2-d_1)$ and odd-character weight-one
series $E_1(\chi_{-3}),E_1(\chi_{-4})$ at rescalings; levels
$\{2,\dots,10,12,13,16,18,20,25\}$; exponent caps $\Sigma|e|\le26$
(pass~1) and $\le60$ (pass~2) with the stated candidate-count caps; and
both $(-q)$-twists of every pair. Two hits are \emph{equivalent} if they
agree under $t\mapsto\lambda t$ and $t\mapsto-t$, detected by the
invariants $a_k/a_1^k$; hit classes are these equivalence classes.
\end{definition}

\begin{census}\label{cen:factory}
\VerifiedR{$\sim5\times10^4$ pairs; all $91$ hit classes re-derived
exactly at series order $46$; \texttt{eps55}}\quad
Within the universe of Definition~\ref{def:universe}: the reverse pipeline
reproduces the exact recurrences of all ten identified families from
their $(t,F)$ data alone ($10/10$); and every $3$-term hit with integral
solution is either one of the fifteen or a classical degenerate
(hypergeometric $c=0$, interleaved/2-step $a=b=0$, or the double-root
case $a^2-4c=0$). \emph{No sixteenth sporadic pair exists in this
universe.} The four uncovered directions (generalized-eta/$\Gamma_1$
banks, Fricke hauptmoduls, order-$4$ territory, levels $>25$) are stated
in \texttt{work/MODULAR\_FACTORY.md}.
\end{census}

%======================================================================
\section{Congruence laws}\label{sec:cong}

\begin{verification}[weight dichotomy]\label{ver:asd}
\VerifiedR{$p\le23$ per family; \texttt{eps51}}\quad
For all nine $r=3$-adjacent (modular weight $2$) families,
$A(p)\equiv b_p\pmod p$ with $b_p=[q^p]F$, at every tested prime ---
including the unidentified Cooper three and the $\chi$-twisted $\zeta$,
$\eta$.
\end{verification}

\begin{proposition}[nebentypus compatibility; a uniform weight-one law]
\label{prop:wt1}
\VerifiedR{fitted on $p\le23$; \emph{passes all sixteen held-out primes}
$29\le p\le97$ for each of the six families below --- the five displayed
plus $\mathbf D$ with its separate rule; \texttt{eps57\_heldout.py}}\quad
For the weight-one families,
\[
A(p)\;\equiv\;b_p\;-\;b\cdot\chi_N(p)\pmod p
\qquad(\mathbf A,\mathbf B,\mathbf C,\mathbf E,\mathbf F),
\]
with $b$ the recurrence's constant parameter $(2,3,3,4,6)$ and $\chi_N$
the quadratic character of the \emph{level} ($\chi_{-3}$ for levels
$6,12,36$ --- including Franel, whose family character is trivial;
$\chi_{-4}$ for level $8$); and for $\mathbf D$,
$A(p)-b_p\equiv2(p\bmod5)-5\pmod p$. For $\zeta$, the programme's twisted
Lucas law holds with twist exponent $e=1$ equal to the nebentypus value
$\chi_{-3}(p)$ at the discriminating primes.
\end{proposition}

\begin{conjecture}[ASD mechanism]\label{conj:asd}
The laws above are the mod-$p$ shadows of Atkin--Swinnerton-Dyer
congruence structure for the identified forms (unit-root/formal-group
statements); in particular the weight-one law should follow from the
Eisenstein structure of the identified weight-one forms. Establishing
this for $\zeta$ --- where Proposition~\ref{prop:zetaid} supplies the
form as a theorem --- is the recommended first target.
\end{conjecture}

%======================================================================
\section{Conjectural tier: mirror-type signatures of the $\zeta(5)$
operator, and a coaction proposal}\label{sec:mirror}

Everything in this section is either an exact local computation about the
Brown--Zudilin operator or an explicitly labelled conjecture; no geometric
origin, Hodge-theoretic, or motivic claim is made.

\subsection{Local signatures}
\begin{proposition}\label{prop:sig}
\Proved\ (exact indicial computation; \texttt{eps50})\quad
The order-$9$ generating operator of the $\Lbz$ recurrence has $t=0$
exponents $0^{(5)}$, $\tfrac12$, and the three roots of an irrational
cubic: a five-dimensional unipotent local block.
\end{proposition}

\begin{verification}\label{ver:mirrorint}
\VerifiedR{$q^{26}$ (mirror map), $q^{32}$ (Yukawa datum), scale one;
\texttt{eps50}, \texttt{eps53}}\quad
The block's mirror map and structure series are integral:
$t(q)=q-94q^2-591q^3-\cdots\in\Z[[q]]$ to the stated order;
$F(q)=1+21q+1015q^2+\cdots\in\Z[[q]]$; and, in the exact $q$-normal form
of the block (defined below), the Yukawa-type datum
$K(q)=\thq^2(\hat g_2/F)=87q+33895q^2+\cdots\in\Z[[q]]$, with
$K_1=87=4P_1$.

\emph{Withdrawn as evidence:} earlier drafts listed the root-integrality
observations $Q(4t)^{1/2},Q(8t)^{1/4}\in\Z[[t]]$ here. These are
\emph{universal} for integer series ($\sqrt{f(4t)}\in\Z[[t]]$ for any
$f\in1+t\Z[[t]]$; see \texttt{papers\_out/half\_apery}, Lemma) and carry
no arithmetic information about this family; the correction is logged in
\texttt{work/Z5\_MODULARITY\_PROBE.md}.
\end{verification}

\begin{definition}[$q$-normal form and the datum $K$]\label{def:normal}
In the nome coordinate of the block, conjugation by $F$ brings the
five-dimensional local system to the form
\[
\thq^{5}\;+\;K_3(q)\,\thq^{3}\;+\;K_2(q)\,\thq^{2}\;+\;K_1(q)\,\thq
\;+\;K_0(q),
\]
with $\hat g_1\equiv0$ (the mirror-map normalisation, verified exactly)
and $\hat g_j$ the successive Frobenius tails; $K:=\thq^2(\hat g_2/F)$ is
the first nontrivial structure series. $K\not\equiv0$ is precisely the
obstruction to the bare inversion of Theorem~\ref{thm:inversion} at
$r=5$: the variation-of-parameters inverse then involves $K$ explicitly.
The full tower $\hat g_3,\hat g_4$ and the invariance of $K$ under the
remaining gauge freedom are \Open\ (referee report \S10); until they are
supplied, $K$'s status is: an exactly computed series in a fixed,
scripted normalisation (\texttt{eps53}).
\end{definition}

\begin{verification}[sharp negatives]\label{ver:mirrornegs}
\ExcludedR{as stated; \texttt{eps53}}\quad
(i) The naive Dwork congruence $F(q)\equiv F_{<p}(q)F(q^p)\bmod p$ fails
at $p=5,7$ (three Frobenius slopes; the slope-refined form is the open
target). (ii) The naive elliptic-style row avatars
$\thq^3(\hat P/Q)$, $\thq^5(P/Q)$ have unboundedly growing denominators
--- calibrated against the same pipeline recovering Beukers' weight-$4$
Eisenstein combination $6E^{(1)}-168E^{(2)}+378E^{(3)}-216E^{(6)}$ exactly
for Ap\'ery ($q^{32}$). (iii) No order-$5$ right factor of $t$-degree
$\le20$ splits the block off globally (one $22$-bit prime, series order
$64$).
\end{verification}

\begin{conjecture}[mirror form of the bridge]\label{conj:mirrorbridge}
In the block's $q$-normal-form coordinates the companion rows
$\hat P_n,P_n$ are coefficient sequences of Eichler-type objects built
from $F$, $K$, and the higher structure series, and the programme's open
compact-formula identity for $P_n$ is the $t$-shadow of a $q$-side
identity in those data.
\end{conjecture}

\subsection{A coaction proposal, in one display}
\begin{conjecture}\label{conj:coaction}
For the motivic period $I^{\mathfrak m}$ of a cellular family with
coaction $\Delta I^{\mathfrak m}=\sum_rI^{\mathrm{dR}}_r\otimes
I^{\mathfrak m}_r$, the $r$-th $\varepsilon$-jet of the family computes
$I^{\mathrm{dR}}_r$; in particular companions are first coaction
components, and the twist behaviour of
Proposition~\ref{prop:wt1} is the Galois equivariance of $\Delta$.
\end{conjecture}
No comparison morphism is constructed here; the conjecture is recorded
because it correctly anticipated the one-step (unipotent) shape of the
first-variation anchor equation of the companion paper, and it makes the
falsifiable prediction that second variations anchor in
$\operatorname{span}\{Q,\dot Q\}$ only.

%======================================================================
\section{Open problems}\label{sec:open}

\begin{enumerate}[leftmargin=2em,label=\textbf{P\arabic*.}]
\item Complete $\delta$'s parametrisation proof (the remaining
Atkin--Lehner degree-one hypothesis), and prove the remaining eleven
identifications by the Ligozat--Sturm route of
Proposition~\ref{prop:zetaid}.
\item Prove the weight-one law of Proposition~\ref{prop:wt1} from the
Eisenstein structure of the identified forms; then the ASD mechanism
(Conjecture~\ref{conj:asd}) for $\zeta$.
\item Certify the second $\zeta$ annihilation law
(Proposition~\ref{prop:ann}); mine the quotient alphabet.
\item Supply the $\hat g_3,\hat g_4$ tower and the gauge-invariance of
$K$ (Definition~\ref{def:normal}); run the slope-refined Dwork test;
derive the $K$-corrected inversion at $r=5$ and test
Conjecture~\ref{conj:mirrorbridge} on the \emph{proved} middle row first.
\item Extend the census (Definition~\ref{def:universe}) in its four
uncovered directions; any of them could still contain a sixteenth pair.
\item Identify Cooper's three parametrizations (Atkin--Lehner
hauptmoduls).
\end{enumerate}

%======================================================================
\section{Provenance and data}\label{sec:provenance}

Scripts in \texttt{work/z5eps/} with logs: \texttt{eps41--eps44}
(negatives), \texttt{eps45} (the proved annihilation law),
\texttt{eps48--eps53} (nome instrument, identifications, companions,
mirror probes), \texttt{eps54} (the discriminating jet experiments),
\texttt{eps55} (census), \texttt{eps56} (half-Ap\'ery),
\texttt{eps57\_sym2\_all.py} and \texttt{eps57\_heldout.py} (the
symmetric-square identities, the $\zeta$ Sturm proof, held-out primes,
Veronese saturation). Memos as in the header. All exact arithmetic is
\texttt{Fraction}-based; every modular negative is certified at two
unrelated primes (Lemma~\ref{lem:modrat}); every positive control is
stated alongside the negatives it guards. The referee report and the
point-by-point response --- including a standing log against silent
evidential upgrades --- are at
\texttt{work/SOL\_REVIEW\_MODULAR\_ANCHORS.md}.

\begin{thebibliography}{99}
\bibitem{Ap79} R.~Ap\'ery, \emph{Irrationalit\'e de $\zeta(2)$ et
$\zeta(3)$}, Ast\'erisque 61 (1979).
\bibitem{Be87} F.~Beukers, \emph{Irrationality proofs using modular
forms}, Ast\'erisque 147--148 (1987).
\bibitem{ASD} A.~O.~L.~Atkin and H.~P.~F.~Swinnerton-Dyer, \emph{Modular
forms on noncongruence subgroups}, Proc.\ Symp.\ Pure Math.\ 19 (1971).
\bibitem{Za09} D.~Zagier, \emph{Integral solutions of Ap\'ery-like
recurrence equations}, CRM Proc.\ Lecture Notes 47 (2009).
\bibitem{AZ06} G.~Almkvist and W.~Zudilin, \emph{Differential equations,
mirror maps and zeta values}, AMS/IP Stud.\ Adv.\ Math.\ 38 (2006).
\bibitem{Co12} S.~Cooper, \emph{Sporadic sequences, modular forms and new
series for $1/\pi$}, Ramanujan J.\ 29 (2012).
\bibitem{BB91} J.~M.~Borwein and P.~B.~Borwein, \emph{A cubic counterpart
of Jacobi's identity and the AGM}, Trans.\ AMS 323 (1991).
\bibitem{Li75} G.~Ligozat, \emph{Courbes modulaires de genre 1},
Bull.\ SMF M\'em.\ 43 (1975).
\bibitem{Stu87} J.~Sturm, \emph{On the congruence of modular forms},
Springer LNM 1240 (1987).
\bibitem{BZ22} F.~Brown and W.~Zudilin, \emph{On cellular rational
approximations to $\zeta(5)$}, 2022--2026.
\bibitem{Br14} F.~Brown, \emph{Irrationality proofs for zeta values,
moduli spaces and dinner parties}, Mosc.\ J.\ Comb.\ Number Theory (2016).
\bibitem{Br26} F.~Brown, \emph{Mellin transforms, the transfinite
diameter, and rational approximation}, preprint, 2026.
\bibitem{Gol07} V.~Golyshev, \emph{Classification problems and mirror
duality}, LMS Lecture Note Ser.\ 338 (2007).
\bibitem{Gor21} O.~Gorodetsky, \emph{New representations for all sporadic
Ap\'ery-like sequences, with applications to congruences}, 2021.
\bibitem{MS16} A.~Malik and A.~Straub, \emph{Divisibility properties of
sporadic Ap\'ery-like numbers}, Res.\ Number Theory 2 (2016).
\bibitem{Sporadics} River's programme, \emph{The fifteen sporadic
Ap\'ery-like pairs}, working paper (\texttt{papers\_out/sporadics}).
\bibitem{VarAnchor} River's programme, \emph{Variational anchors for
Ap\'ery-like families}, working paper
(\texttt{papers\_out/variational\_anchor}).
\bibitem{HalfAp} River's programme, \emph{The half-Ap\'ery numbers},
working paper (\texttt{papers\_out/half\_apery}).
\bibitem{ExpAp} River's programme, \emph{The companion by the nome},
expository paper (\texttt{papers\_out/expository\_apery}).
\end{thebibliography}

\end{document}
